\documentclass[11pt,a4paper]{article}
\usepackage[margin=2.5cm]{geometry}
\usepackage{amsmath,amssymb}
\usepackage[hidelinks]{hyperref}
\newcommand{\doi}[1]{\href{https://doi.org/#1}{doi:#1}}
\setlength{\emergencystretch}{4em}

% --- notation shared with traceable_energetics.tex ----------------------------------
\newcommand{\vt}{\vartheta}
\newcommand{\code}[1]{\texttt{#1}}
% this note
\newcommand{\Mtot}{M_{\rm tot}}
\newcommand{\Ptop}{P_{\rm top}}
\newcommand{\zs}{z_{*}}

\title{The Lorenz reference state of available energy\\ in AN2D and FC2D: physics and algorithms}
\author{Notes on the \code{TraceableEnergetics} code}
\date{}

\begin{document}
\maketitle
\tableofcontents

\bigskip
\noindent\textbf{Scope.}
This note describes how \code{TraceableEnergetics} computes the Lorenz reference state of a snapshot of AN2D or FC2D, the
available energy of each parcel, the tendency of the reference state and the budgets of the available energy
(\code{src/reference\_state/}, with the numerical routines in \code{src/numerics/}). The available energy and its budget are
derived in the energy-budget notes (\code{energy\_budget/energy.tex}, ``the energy notes''); notation follows
\code{models/traceable\_energetics.tex} (``the main note''). Appendix~\ref{app:future} lists what was specified or considered
but is not implemented.

\paragraph{Notation.} $\vt$ is the conserved variable (\code{consvar}), $q$ the solute fraction, $h(P,\vt,q)$ the specific
enthalpy with $\partial h/\partial P=v$ the specific volume, and $\Phi=g(z-z_0)$. The box is $[0,L_x]\times[0,L_z]$ with
$M_x\times M_z$ interior cells of size $\Delta x\times\Delta z$; masses are per unit length along $y$. AN2D has the reference
profiles $\rho_R(z)$, $P_R(z)$ with $\partial_zP_R=-g\rho_R$. Reference-state quantities carry a subscript $*$: the reference
pressure $P_*(z,t)$, the reference specific volume $v_*(z,t)$ (density $\rho_*=1/v_*$), the reference height $\zs$ of a
parcel, and in AN2D the reduced reference pressure $\pi_*=(P_*-P_R)/\rho_R$.

% =====================================================================================
\section{Available energy and the reference state}\label{sec:theory}
The static energy of a parcel at height $z$ and pressure $P$ is
\begin{equation}
  \Sigma(\vt,q,z,P)=\begin{cases}
    h(P,\vt,q)+\Phi(z) & \text{FC2D},\\[2pt]
    h(P_R(z),\vt,q)+\dfrac{P-P_R(z)}{\rho_R(z)}+\Phi(z) & \text{AN2D},
  \end{cases}
  \qquad \Sigma_P=\begin{cases}v & \text{FC2D},\\ 1/\rho_R & \text{AN2D}.\end{cases}
  \label{eq:Sigma}
\end{equation}
Given a reference pressure $P_*(z,t)$, the available energy (AE) of a parcel is
\begin{equation}
  \Pi=\Sigma(\vt,q,z,P)-\Sigma\left(\vt,q,\zs,P_*(\zs)\right)-\left(P-P_*(z)\right)\Sigma_P(\vt,q,z,P)=\Pi_1+\Pi_2,
  \label{eq:A}
\end{equation}
the sum of the available compressible energy (ACE), the compression from $P_*(z)$ to $P$ at fixed height,
\begin{equation}
  \Pi_1=\Sigma(\vt,q,z,P)-\Sigma(\vt,q,z,P_*(z))-\left(P-P_*(z)\right)\Sigma_P(\vt,q,z,P)\ \ge0,
\end{equation}
zero in AN2D, where $\Sigma$ is linear in $P$, and nonnegative in FC2D, where $h$ is concave in $P$ ($\partial_Pv<0$); and the
available potential energy (APE), the displacement from $\zs$ to $z$ through the reference state,
\begin{equation}
  \Pi_2=S(z)-S(\zs),\qquad S(z)=\Sigma\left(\vt,q,z,P_*(z)\right),
  \label{eq:Pi2}
\end{equation}
with $S$ the static energy of the parcel along the reference profile. The reference profile defines $v_*$ by
\begin{equation}
  v_*\,\partial_zP_*=-g\quad\text{(FC2D)},\qquad \partial_z\pi_*=g\left(\rho_Rv_*-1\right)\quad\text{(AN2D)}.
  \label{eq:vstar}
\end{equation}
The energy notes require of the reference state:
\begin{enumerate}
  \item \textbf{Level of neutral buoyancy (LNB).} $S'(\zs)=0$, i.e.\ $v(P_*(\zs),\vt,q)=v_*(\zs)$ (FC2D) or
        $v(P_R(\zs),\vt,q)=v_*(\zs)$ (AN2D), with $N_*^2=S''(\zs)\ge0$;
  \item \textbf{Volume constraint.} The reference state fills the box level by level: for every $f$,
        \begin{equation*}
          \int f(\zs)\,\Sigma_P(\vt,q,\zs,P_*(\zs))\,\mathrm{d}m=\int f(z)\,\mathrm{d}V .
        \end{equation*}
\end{enumerate}
These are the stationarity conditions of $\int\Pi\,\mathrm{d}m$ with respect to $\zs$ and $P_*$.

\paragraph{Potential energies.} The potential energy of a state is that of the energetics of the model, the integral of
$\rho\,(h-P/\rho+\Phi)$, the internal plus gravitational energy (per unit length along $y$):
\begin{equation}
  E_P=\begin{cases}
    \displaystyle\int\left[\rho\left(h(P,\vt,q)+\Phi\right)-P\right]\mathrm{d}A & \text{FC2D},\\[8pt]
    \displaystyle\int\left[\rho_R\left(h(P_R,\vt,q)+\Phi\right)-P_R\right]\mathrm{d}A & \text{AN2D},
  \end{cases}
  \label{eq:EP}
\end{equation}
where AN2D evaluates the thermodynamics at the reference pressure $P_R(z)$, as its energy does. The \textbf{reference
energy} $E_*$ is the same integral for the reference state: every fluid element of the actual state, of mass
$\mathrm{d}m=\rho\,\mathrm{d}A$ and conserved $(\vt,q)$, moved adiabatically to its reference height $\zs(\vt,q)$, at rest
under the reference pressure. Written over the mass, with the reference state filling the box (the volume constraint),
\begin{equation}
  E_*=\begin{cases}
    \displaystyle\int\left[h\left(P_*(\zs),\vt,q\right)+\Phi(\zs)\right]\mathrm{d}m-L_x\int_0^{L_z}P_*(z)\,\mathrm{d}z & \text{FC2D},\\[10pt]
    \displaystyle\int\left[h\left(P_R(\zs),\vt,q\right)+\Phi(\zs)\right]\mathrm{d}m-L_x\int_0^{L_z}P_R(z)\,\mathrm{d}z & \text{AN2D},
  \end{cases}
  \label{eq:Estar}
\end{equation}
the integrals in $\mathrm{d}m$ over the fluid (each element at its own $\zs$), the pressure integrals over the box. In FC2D,
$P_*$ is hydrostatic with the density $\rho_*$ of the rearranged fluid, $\partial_zP_*=-g\rho_*$; integrating by parts,
$L_x\int_0^{L_z}P_*\,\mathrm{d}z=\Ptop L_xL_z+\int\Phi(\zs)\,\mathrm{d}m+gz_0\Mtot$, so
\begin{equation}
  E_*=\int h\left(P_*(\zs),\vt,q\right)\mathrm{d}m-\Ptop L_xL_z-gz_0\Mtot\qquad\text{(FC2D)},
\end{equation}
with $\Ptop=P_*(L_z)$. In AN2D the mass per unit height of the reference state is that of the model, $L_x\rho_R(z)$, and the
pressure term is the same in $E_P$ and $E_*$.

The \textbf{Lorenz reference state} is the rearrangement (the map of the fluid elements to their heights $\zs$, with the
volume constraint, and in FC2D with $P_*$ hydrostatic with it) that minimises $E_*$. Then
$E_P-E_*\ge0$ (exactly in AN2D, up to an $O(\Delta z^4)$ discretisation difference in FC2D, \S\ref{sec:slots}) is the
potential energy that any adiabatic rearrangement can release, the available potential energy of the whole box in the
sense of Lorenz. It is a number for the box, not a sum of parcel energies. When the volume constraint holds,
\begin{equation}
  \int\Pi\,\mathrm{d}m=E_P-E_* ,
  \label{eq:PiE}
\end{equation}
since, from \eqref{eq:A}, the terms of $\Pi$ at the actual state give $E_P$ plus $\int P\,\mathrm{d}V$ (FC2D; $\int P_R\,\mathrm{d}A$
plus $\int(P-P_R)\,\mathrm{d}A$ in AN2D), those at the reference state $E_*$ plus $L_x\int P_*\,\mathrm{d}z$, and the
pressure term $-\int(P-P_*(z))\,\Sigma_P\,\mathrm{d}m=-\int(P-P_*(z))\,\mathrm{d}V$ cancels the rest. For any other profile
with the global minimisers $\zs$, $\int\Pi\,\mathrm{d}m\ge E_P-E_*$ of the Lorenz state (\S\ref{sec:dual}). The diagnostic
\code{background\_potential\_energy} is $E_P-\int\Pi\,\mathrm{d}m$, which is $E_*$ when \eqref{eq:PiE} holds. The discrete
$E_P$ and $E_*$ are given in \S\ref{sec:slots}.

The code takes
$\zs$ to be the \emph{global} minimiser of $S$ on $[0,L_z]$ (\S\ref{sec:lnb}): it maximises the parcel's $\Pi$, depends on
$(\vt,q,t)$ only, gives $\Pi_2\ge0$, and makes the terms in $\partial_t\zs$ drop out of the budgets.

Discretely (\S\ref{sec:problem}) the Lorenz reference state is a placement of the parcels into the levels of a column, whose constraints are the
discrete volume constraint and whose dual potentials are the discrete $P_*$.

% =====================================================================================
\section{The rearrangement problem}\label{sec:problem}
\subsection{Parcels}
Each interior cell is a parcel $n=1,\dots,N$, $N=M_xM_z$, numbered row by row from the bottom. It keeps its mass and
conserved properties,
\begin{equation}
  m_n=\rho\,\Delta x\Delta z,\qquad \vt_n=\frac{(\rho\vt)_{ij}}{\rho},\qquad q_n=\frac{(\rho q)_{ij}}{\rho},
\end{equation}
with $\rho=\rho_R(z_j)$ in AN2D and the cell density in FC2D. It also has its own pressure $p_n$ ($P_R(z_n)$ in AN2D; in FC2D
the pressure $p(v_n,\vt_n,q_n)$ of the equation of state at $v_n=1/\rho$) and height $z_n$, the centre of its row.

\subsection{Slots and costs}\label{sec:slots}
The reference column has $K=M_z$ \emph{slots}, the model rows, numbered from the top, of masses $M_k$ ($\sum_kM_k=\Mtot$). A
placement $x_{nk}\ge0$ is the mass of parcel $n$ in slot $k$; the Lorenz state is the least-cost placement,
\begin{equation}
  \min_{x\ge0}\ \sum_{n,k}x_{nk}\,c_{nk}\qquad\text{subject to}\qquad \sum_kx_{nk}=m_n\ \ \forall n,\qquad \sum_nx_{nk}=M_k\ \ \forall k ,
  \label{eq:tp}
\end{equation}
the transportation problem (Hitchcock 1941). It generalises the linear assignment problem of Hieronymus \& Nycander (2015)
(equal masses, $K=N$) and the minimum-cost flow problem of Su \& Ingersoll (2016) (equal masses, $K\ll N$ levels): the
parcels have unequal masses here, and may be split between slots. The costs are energies per unit mass measured from the
parcel's own state, so that a constant per parcel, which changes neither the placement nor the slot duals, cancels:
\begin{itemize}
  \item \textbf{AN2D.} The slots are the rows at the fixed pressures $P_k=P_R(z_k)$, of masses $M_k=M_x\rho_R(z_k)\Delta x\Delta z$.
        These are the states anelastic flow can reach ($\nabla\cdot(\rho_R\mathbf u)=0$ preserves $\rho_R\,\mathrm{d}A$). The
        cost is
        \begin{equation}
          c_{nk}=h(P_R(z_k),\vt_n,q_n)-h(P_R(z_n),\vt_n,q_n)+g\,(z_k-z_n),
          \label{eq:an-cost}
        \end{equation}
        zero in the parcel's own row. The discrete energies \eqref{eq:EP} of the state and of a placement are
        \begin{equation}
          E_P=\sum_n\left[m_n\left(h(P_R(z_n),\vt_n,q_n)+\Phi(z_n)\right)-P_R(z_n)\,\Delta x\Delta z\right],
        \end{equation}
        \begin{equation}
          E_*=\sum_{nk}x_{nk}\left(h(P_R(z_k),\vt_n,q_n)+\Phi(z_k)\right)-\sum_kM_x\,P_R(z_k)\,\Delta x\Delta z ,
        \end{equation}
        $E_P$ being the integrated potential energy of the energetics. The rows keep their masses, so the pressure terms
        are equal, and $E_P-E_*=-\sum_{nk}x_{nk}c_{nk}$ exactly. The actual state is the feasible placement ``each
        parcel in its own row'', so $E_*\le E_P$ to round-off for the optimal placement, with equality for a state at rest
        and stably stratified.
  \item \textbf{FC2D.} The slots are the rows with the masses $M_k$ of the rows of the state, so that a state at rest is its
        own reference state. The column is hydrostatic in mass, and slot $k$ is at the pressure
        \begin{equation}
          P_k=\Ptop+\frac{g}{L_x}\Big(\sum_{l<k}M_l+\tfrac12M_k\Big),\qquad c_{nk}=h(P_k,\vt_n,q_n)-h(p_n,\vt_n,q_n),
          \label{eq:fc-cost}
        \end{equation}
        with the top pressure $\Ptop$ set by the condition that the column fills the box (\S\ref{sec:fc}). Slot $k$ has the
        volume $V_k=\sum_nx_{nk}v(P_k,\vt_n,q_n)$ and the thickness $V_k/L_x$, stacked from the bottom. Integrating
        $\rho gz$ by parts in the column, $\sum_kM_kgz_k-\sum_kP_kV_k=-\Ptop L_xL_z$. The discrete energies \eqref{eq:EP}
        of the state and of the column are then
        \begin{equation}
          E_P=\sum_n\left[m_n\left(h(p_n,\vt_n,q_n)+\Phi(z_n)\right)-p_n\,\Delta x\Delta z\right],
        \end{equation}
        \begin{equation}
          \begin{aligned}
            E_*&=\sum_{nk}x_{nk}\left(h(P_k,\vt_n,q_n)+\Phi(z_k)\right)-\sum_kP_kV_k\\
               &=\sum_{nk}x_{nk}\,h(P_k,\vt_n,q_n)-gz_0\Mtot-\Ptop L_xL_z ,
          \end{aligned}
        \end{equation}
        the geopotential of the column independent of the placement at fixed $\Ptop$, and
        \begin{equation}
          E_P-E_*=\sum_n\left[m_ngz_n-(p_n-\Ptop)\Delta x\Delta z\right]-\sum_{nk}x_{nk}c_{nk}.
        \end{equation}
        For each placement, $\Ptop$ is the pressure at which its column fills the box, and the Lorenz state is the
        placement of least $E_*$ (the saddle point of \S\ref{sec:fc}).
        For a state at rest, $E_P-E_*$ is the difference between the model's discrete hydrostatic balance and that of the
        column, $O(\Delta z^4)$ (about $10^{-18}\Mtot gL_z$ at $M_z=8$).
\end{itemize}
The enthalpy differences are integrals of $v$ by three-point Gauss--Legendre quadrature,
$h(P+\Delta P)-h(P)=\int v\,\mathrm{d}P$ (exact for the equation of state of the main note, linear in $P$). This keeps the
relative precision of the float type: the difference of two values of $h\approx10^6$\,J\,kg$^{-1}$ would carry
$2\times10^{-10}$\,J\,kg$^{-1}$ of round-off in Float64, far above the tolerance of the certificate below.

\subsection{Duals and the certificate}\label{sec:dual}
The dual of \eqref{eq:tp} has a potential $u_n$ per parcel and $\psi_k$ per slot. A placement is optimal if and only if
there are duals with
\begin{equation}
  c_{nk}-u_n-\psi_k\ge0\ \ \forall n,k,\qquad\text{with equality wherever }x_{nk}>0
  \label{eq:dual}
\end{equation}
(weak duality and complementary slackness; e.g.\ Bertsimas \& Tsitsiklis 1997, ch.~4). The code checks the first condition
for all $NK$ pairs after every placement, the second holding by construction: \code{certificate\_violations} counts the
parcels whose least reduced cost $\min_k(c_{nk}-u_n-\psi_k)$ is below $-16(K+1)\,\epsilon\,gL_z$, the round-off of duals that
accumulate $K$ cost differences of size up to $gL_z$ ($\epsilon$ the machine epsilon of the float type). This pass costs
$NK$ evaluations of the cost, threaded and vectorised over the parcels of a row, and is also the pricing of the exact method
(\S\ref{sec:exact}).

The duals are the discrete reference pressure. In AN2D, $S(z_k)=c_{nk}+\pi_*(z_k)$ up to a constant per parcel, so a parcel's
optimal slots minimise $c_{nk}-\psi_k$ and $\pi_{*,k}=-\psi_k$ up to a constant; in FC2D the duals give $-\Phi(z_k)$. Dual
feasibility with complementary slackness is the discrete LNB condition in global form, and $\psi$ is the multiplier of the
slot constraints, the discrete volume constraint. By weak duality, for \emph{any} reference profile, $\int\Pi\,\mathrm{d}m$ with
the global minimiser $\zs$ is at least $E_P-E_*$ of the Lorenz state (cf.\ Tailleux 2013), with equality in AN2D for the optimal
duals with $\zs$ restricted to the row centres. The code uses the duals only to certify the placement; the profile comes from
the slot volumes (\S\ref{sec:profile}).

\subsection{When the sort is exact}\label{sec:why}
The equation of state of the main note (Appendix A there) is
\begin{equation}
  v = v_0\left[1+\alpha_T\left(1+\gamma(P-P_0)\right)(\vt-T_0)-\alpha_q(q-q_0)-\alpha_P(P-P_0)+\tfrac12\alpha_{TT}(\vt-T_0)^2\right].
  \label{eq:eos}
\end{equation}
If $v$ is additively separable, $v=f(\vt,q)+w(P)$, then $h(P,\vt,q)=h(P_0,\vt,q)+f(\vt,q)(P-P_0)+\int_{P_0}^Pw$, and only
$\sum f_nP_k$ depends on the placement. With the parcels ordered by $f$ and the slots by pressure the costs have the Monge
property, for which filling the slots in order (the north-west corner rule) is optimal, with unequal masses and split
parcels (Hoffman 1963). In FC2D the slot pressures depend on $\Ptop$ but their order does not. So \textbf{the sort by specific
volume is the exact Lorenz state iff $\gamma=0$}: compressibility, haline contraction and cabbeling keep $v$ separable,
thermobaricity does not. With $\gamma\ne0$ the order of two parcels can depend on the pressure at which they are compared,
and the sort is approximate. In the Boussinesq limit (uniform $\rho_R$, equal masses) and with $\gamma=0$ it is the sort of
Winters et al.\ (1995).

% =====================================================================================
\section{Algorithms of the placement}\label{sec:algorithms}
The method is a type parameter of the reference state: \code{SortedRearrangement()} (the default of the constructors) or
\code{ExactRearrangement()} (used by the examples). Both start from the sort, so the two can be compared on the same state.

\subsection{Sorted rearrangement}\label{sec:sort}
\begin{enumerate}
  \item The sort key of each parcel is its specific volume at a key pressure, $v(P_{\rm key},\vt_n,q_n)$: $P_R(L_z/2)$ in AN2D,
        the mean pressure of the parcels in FC2D. With $\gamma=0$ every $P_{\rm key}$ gives the same order.
  \item The parcels, lightest first, fill the slots from the top, a parcel split across a slot boundary where needed. The
        entries of the placement form a staircase.
  \item FC2D: $\Ptop$ for this placement (\S\ref{sec:fc}).
  \item The duals along the staircase: $\psi_1=0$, and across the boundary between slots $k$ and $k+1$ the step of $\psi$ is
        $c_{a,k+1}-c_{a,k}$ for the parcel $a$ split there; at a degenerate boundary (no split parcel), the middle of the
        interval allowed by the last parcel of slot $k$ and the first of slot $k+1$, a reproducible choice. Then
        $u_n=c_{nk}-\psi_k$ on any slot of parcel $n$.
  \item One pass over the pairs: the certificate.
\end{enumerate}
Cost: $O(N\log N)$ for the sort, $O(NK)$ for the certificate.

\subsection{Exact rearrangement: the transportation simplex}\label{sec:exact}
A transportation simplex (Dantzig 1951), in the network form of Ahuja, Magnanti \& Orlin (1993, ch.~11), written in Julia
(\code{src/numerics/transportation.jl}). The network has $N+K$ nodes, the parcels and the slots, with an arc from each parcel
to each slot carrying $x_{nk}$; a basis is a spanning tree of $N+K-1$ arcs, some of which may carry no mass.
\begin{enumerate}
  \item \textbf{Start.} The sorted placement, completed into a spanning tree: at a degenerate boundary between slots $k$ and
        $k+1$, the first parcel of slot $k+1$ gets an entry of zero mass in slot $k$.
  \item \textbf{Duals.} At the start of each pass, $\psi_1=0$ and $u_n+\psi_k=c_{nk}$ on the arcs of the tree, by a
        breadth-first walk from the first slot, in $O(N+K)$ with the cost of each arc kept.
  \item \textbf{Pricing.} A pass over all $NK$ pairs (the certificate pass of \S\ref{sec:dual}) gives each parcel's least
        reduced cost and its slot. The parcels below $-16(K+1)\epsilon gL_z$ are the candidates, sorted once, most negative first
        (the lower index first among equals). Each is priced again with the current duals, one cost evaluation, and enters
        the tree if still negative. Taking the most favourable arc of each node as a candidate is the candidate-queue pricing
        of large network simplex codes (Bradley, Brown \& Graves 1977; see also Grigoriadis 1986). The passes repeat until
        one finds no candidate: that pass is the certificate.
  \item \textbf{Pivot.} The entering pair $(n,k)$ closes a cycle in the tree, from $n$ up to the apex (the deepest common
        ancestor of $n$ and $k$) and down to $k$. A mass $\theta$ moves round it, $+\theta$ on the arcs oriented as the entering
        one and $-\theta$ on the others; $\theta$ is the least mass of these, so masses stay exactly nonnegative in floating
        point. The leaving arc is the last arc that blocks $\theta$ met going round the cycle from the apex in the direction
        of the entering arc.
  \item \textbf{Update of the tree.} Only the subtree that hung from the leaving arc changes: it now hangs from the entering
        arc, rerooted at the end $q$ of that arc that it holds, and its depths and duals change, in a time of the size of
        the subtree rather than $O(N+K)$ (as with the threaded indices of Glover, Klingman \& Stutz 1974; here with lists of
        children). With $r$ the reduced cost of the entering pair, the duals of the subtree shift by $u\to u+r$,
        $\psi\to\psi-r$ if $q$ is the parcel, and $u\to u-r$, $\psi\to\psi+r$ if $q$ is the slot. Step 2 recomputes the duals
        from the costs at each pass, so the certificate carries the round-off of the costs only.
\end{enumerate}
The tree is rooted at the first slot and kept \emph{strongly feasible} (Cunningham 1976): every arc of zero mass goes from a
parcel up to its parent slot, which the completed staircase satisfies and the leaving rule preserves. A pivot that moves no
mass, frequent from the degenerate staircase, then still makes progress, and the method cannot cycle. The optimum is global:
\eqref{eq:tp} is a linear programme, and the certificate proves it. The pivots give the decrease of the cost,
$\sum\theta\,|c_{nk}-u_n-\psi_k|$, by which the sort underestimates $E_P-E_*$ (\code{sorted\_energy\_excess}).

The number of passes grows with how far the sort is from the optimum: one where it is optimal, a few where it fails for a few
nearly tied parcels, about $K$ where every parcel is misplaced, as each pass moves a misplaced parcel about one slot nearer its
level. The passes are capped at $10K$ (\code{passes\_per\_slot}), which only guards against round-off; a placement cut short
fails its certificate.

\subsection{FC2D: the top pressure}\label{sec:fc}
For a placement $x$, let $G(x,\Ptop)=\sum x_{nk}c_{nk}(\Ptop)-\Ptop L_xL_z$, the energy of the column up to terms independent
of both. Then $\partial G/\partial\Ptop=\sum x_{nk}v(P_k,\vt_n,q_n)-L_xL_z$ is the volume excess and
$\partial^2G/\partial\Ptop^2=\sum x_{nk}\partial_Pv<0$: $G$ is concave in $\Ptop$, with its maximum where the column fills the
box. The reference state is the saddle point
\begin{equation}
  \min_x\ \max_{\Ptop}\ G(x,\Ptop)=\max_{\Ptop}\ \min_x\ G(x,\Ptop)
  \label{eq:saddle}
\end{equation}
($G$ is linear in $x$ over a compact convex set and concave in $\Ptop$; Sion's minimax theorem), the discrete counterpart of
the statement of the energy notes that $\int\Pi\,\mathrm{d}m$ is a minimum with respect to $P_*$, and of the free-surface
optimal transport of Egan et al.\ (2025).
\begin{itemize}
  \item \textbf{$\Ptop$ at fixed placement} (\code{top\_pressure!}): Newton's method on the volume residual
        $R(\Ptop)=\sum x_{nk}v(P_k)-L_xL_z$, with $R'=\sum x_{nk}\partial_Pv(P_k)$, from the mean over the top row of
        $p_n-g\Delta z/(2v_n)$. It stops when $|R|\le16\epsilon L_xL_z$ or $|R|$ no longer decreases.
  \item \textbf{Sorted method:} the placement does not depend on $\Ptop$ when $\gamma=0$, so $\Ptop$ is computed once.
  \item \textbf{Exact method:} from the sorted placement, $\Ptop$ at fixed placement and the simplex at fixed $\Ptop$ (from the
        previous tree) alternate until a simplex moves no mass, at most 20 times (\code{max\_top\_iterations}); the last
        placement is then only priced, and its certificate fails. At the end both conditions hold: the placement is optimal
        at $\Ptop$, and $\Ptop$ fills the box with it, which is the saddle point.
\end{itemize}

\subsection{Complexity}\label{sec:complexity}
\paragraph{Of the problem.} \eqref{eq:tp} is a transportation problem, an uncapacitated minimum-cost flow problem on
$N+K$ nodes and $NK$ arcs, so it is solvable in strongly polynomial time, in a number of operations bounded by a polynomial
in $N$ and $K$ alone (Tardos 1985). Orlin's (1993) algorithm solves uncapacitated minimum-cost flow on $n$ nodes and $m$ arcs
in $O(n\log n\,(m+n\log n))$, here $O(N^2K\log N)$; Kleinschmidt \& Schannath (1995) give a strongly polynomial algorithm for
the transportation problem itself. With few demand points, $K\ll N$, Tokuyama \& Nakano (1995) solve it in expected
$O(KN)$, which is optimal, when $N=\Omega(K^{4+\varepsilon})$ and the ratio of the largest to the smallest parcel mass is
polynomial in $N$; our grids ($N=2048$, $K=32$) are far from that regime. Without thermobaricity the costs are Monge
(\S\ref{sec:why}) and the problem is a sort, $O(N\log N)$. Equal masses with $K=N$ is linear assignment, $O(N^3)$.

\paragraph{Of the algorithm used.} The primal network simplex with a strongly feasible tree terminates (Cunningham 1976),
but no polynomial bound on the number of its pivots is known for pricing rules such as ours, the most negative reduced cost
of each parcel; with the closely related rule of Dantzig, the most negative reduced cost overall, network simplex methods take
exponentially many pivots on some networks (Zadeh 1973).
Polynomial pivot rules exist (Orlin 1997). The observed cost is (passes)$\,\times NK$ cost evaluations, plus small
per-pivot terms (\S\ref{sec:exact}): one pass where the sort is optimal, a few where it fails for a few nearly tied parcels,
and at most about $K$ in the cases tested (1.0 to $1.3K$ where every parcel is misplaced, as each pass moves a misplaced parcel
about one slot nearer its level), i.e.\ $O(NK^2)$ evaluations. This is an observation on our states, not a theorem. The
certificate proves each result optimal however it was found, so a solver with a guaranteed bound could replace or back the
simplex without changing anything else (Appendix~\ref{app:future}).

% =====================================================================================
\section{The reference profile}\label{sec:profile}
The profile is built from the slot-mean specific volume and compressibility of the parcels of each slot,
\begin{equation}
  v_{*,k}=\frac1{M_k}\sum_nx_{nk}\,v(P_k,\vt_n,q_n),\qquad b_k=\frac1{M_k}\sum_nx_{nk}\,\partial_Pv(P_k,\vt_n,q_n),
\end{equation}
which change continuously with the state, unlike the duals (\S\ref{sec:decisions}). The pressure is interpolated, and $v_*$
derived from the interpolant, so that $\Pi$, $v_*$, $N_*^2$ and the LNB residual come from one profile and $\zs$ is an exact
minimiser of the $S$ in $\Pi$.
\begin{itemize}
  \item \textbf{AN2D.} $P_R$ is the $C^1$ cubic Hermite interpolant through the cell centres with the slopes $-g\rho_R$. $\pi_*$
        is a cubic Hermite interpolant with nodes at the row centres and both walls. At the row centres its slopes are
        $g(\rho_Rv_{*,k}-1)$ and its values follow from $\partial_z\pi_*=-v_*\partial_zP_R-g$ with $v_*$ linear in $P_R$
        between the rows (the trapezoidal rule in the pressure),
        \begin{equation}
          \pi_{*,c}-\pi_{*,c-1}=-\tfrac12(v_{*,c-1}+v_{*,c})\left(P_R(z_c)-P_R(z_{c-1})\right)-g(z_c-z_{c-1}),
        \end{equation}
        exact when $v$ is linear in $P$, so that the static energy of a neutral fluid is flat at the nodes. Between the end
        row $e$ and the wall, $v_*$ follows the compression of the parcels of the end slot, $v_*=v_{*,e}+b_e(P_R-P_R(z_e))$,
        which gives the wall nodes
        $\pi_*=\pi_{*,e}-v_{*,e}\Delta P-\tfrac12b_e\Delta P^2-g(z-z_e)$, $\Delta P=P_R(z)-P_R(z_e)$. (A constant $v_*$ there
        lowered the static energy of these parcels towards the walls by $\tfrac12|b_e|(g\rho_R\Delta z/2)^2$ and made a wall
        the reference height of every parcel of a neutral fluid.) $\pi_*$ is defined up to a constant, which changes neither
        $\zs$ nor $\Pi$; its mean over the mass of the rows is zero. On the model rows, $P_*=P_R+\rho_R\pi_*$ and
        $v_*=(\partial_z\pi_*/g+1)/\rho_R$.
  \item \textbf{FC2D.} The slots are stacked from the bottom with the thicknesses $V_k/L_x$; edge $e$ has the pressure $\Ptop$
        plus $g/L_x$ times the mass above it. $P_*$ is the monotone $C^1$ cubic (PCHIP: Fritsch--Carlson slopes with
        Brodlie's weights, three-point end slopes) through the edge pressures, with the secants $-g/v_{*,k}$, so that the
        slot density $M_k/V_k$ is exactly the mean of $-\partial_zP_*/g$ across the slot. On the model rows,
        $v_*=-g/\partial_zP_*$. The column fills the box to the round-off of $\Ptop$ (\code{volume\_residual}).
\end{itemize}

% =====================================================================================
\section{Reference heights}\label{sec:lnb}
The reference height of a parcel is the global minimiser of its static energy along the profile,
\begin{equation}
  \zs=\arg\min_{0\le z\le L_z}S(z),\qquad
  r(z)=-S'(z)=\begin{cases}
    g\rho_R\left(v(P_R(z),\vt,q)-v_*(z)\right) & \text{AN2D},\\[2pt]
    g\left(\rho_*v(P_*(z),\vt,q)-1\right) & \text{FC2D},
  \end{cases}
  \label{eq:zstar}
\end{equation}
with $r$ the LNB residual, $-(v\,\partial_zP_R+g+\partial_z\pi_*)$ in AN2D and $-(v\,\partial_zP_*+g)$ in FC2D before
\eqref{eq:vstar} is used ($r>0$: the parcel is lighter than the reference state at $z$), and
$N_*^2=S''=\partial_Pv\,(\partial_zP)^2+v\,\partial_z^2P+\partial_z^2\pi_*$ (no $\pi_*$ in FC2D), which depends on the parcel
through $(\vt,q)$ (Tailleux 2013, Appendix A). $S$ is measured from the parcel's own pressure and height, with the enthalpy
differences by quadrature. For each parcel, by the model's manager:
\begin{enumerate}
  \item \textbf{Coarse search}, vectorised over the parcels of a row: $S$ at the levels of the profile, both walls and the
        row centres in AN2D, both walls and the interior slot edges in FC2D, and the lowest (the lowest level among equal
        values).
  \item \textbf{Refinement} next to that level: if $r$ changes sign towards a neighbouring level, its root, found by the
        bracketing method of Alefeld, Potra \& Shi (1995) to within $\sqrt\epsilon\,L_z$ (which leaves $S$ at round-off), is
        kept if $S$ is lower there. A wall is the minimum where $r$ has the outward sign.
  \item \textbf{The parcel's own height} is the reference height if $S$ is lower still there (in FC2D it is not a level, and
        in a well-mixed fluid, where $S$ is flat to round-off, it can be the lowest point found).
\end{enumerate}
A local minimum next to another level is not refined; it can be lower than the one found only if the two differ by less than
the change of $S$ between two levels. Multiple minima arise only through thermobaricity. The global minimiser (rather than
the nearest root of Tailleux 2013, \S3.1) makes $\zs$ independent of the parcel's height, so $\Pi$ does not jump when a parcel
crosses a local maximum of $S$; $S(\zs)$ is continuous in $(\vt,q,t)$ even where $\zs$ jumps between two minima. The number of
parcels at a wall is a diagnostic (\code{wall\_minima}).

For a parcel in a slot that holds a mixture, $\zs$ can differ from its placement by up to a slot thickness, as $v_*$ is a slot
average: $\zs$ satisfies the LNB condition, which the local budgets need, and the placement the volume constraint.

% =====================================================================================
\section{Available energy}\label{sec:ae}
On the cells: $\Pi_2=S(z_n)-S(\zs)\ge0$, with both values measured from the parcel's own state as in the search, so that
$\Pi_2\ge0$ holds exactly; $\Pi_1=0$ in AN2D; in FC2D
\begin{equation}
  \Pi_1=h(p)-h(P_*(z))-(p-P_*(z))\,v(p)=-\int_{P_*(z)}^{p}(s-P_*(z))\,\partial_Pv(s,\vt,q)\,\mathrm{d}s\ \ge0,
\end{equation}
by quadrature, with $P_*(z)$ at the centre of the parcel's row. Their integrals are $\sum_nm_n\Pi_n$; the background potential
energy is $E_P-\sum_nm_n\Pi_n$, with $E_P$ the integrated potential energy of the energetics.

% =====================================================================================
\section{Tendency of the reference state}\label{sec:dtpstar}
The budgets need $\partial_tP_*$ on the model rows and at each $\zs$. The Lorenz state depends only on the distribution of mass
over $(\vt,q)$, which advection does not change, so $P_*$ changes only through the diabatic rates (Winters et al.\ 1995;
Winters \& D'Asaro 1996)
\begin{equation}
  \dot\vt=\frac{\sigma_\vt-\nabla\cdot\mathbf J_\vt}{\rho},\qquad \dot q=-\frac{\nabla\cdot\mathbf J_q}{\rho},\qquad
  \dot v=\partial_\vt v\,\dot\vt+\partial_qv\,\dot q\ \ \text{(at fixed pressure)},
\end{equation}
which the model provides ($\rho=\rho_R$ in AN2D). The tendency is computed from a single snapshot (\code{reference\_tendency!}).

\paragraph{Derivation.} Follow a parcel at its reference level and differentiate the LNB condition. In FC2D, in the mass
coordinate $m$ (the mass above a level), $P_*(m,t)=\Ptop(t)+gm/L_x$, and with $N_*^2=g\rho_*^2(v_*\partial_zv_*+gv_P)$,
\begin{equation}
  \frac{N_*^2}{g\rho_*^2L_x}\,\frac{Dm_*}{Dt}=\partial_tv_*|_m-\dot v-v_P\,\partial_t\Ptop .
\end{equation}
The mass above a level stays the same, so the mass-weighted mean of $Dm_*/Dt$ over the parcels at a level vanishes:
\begin{equation}
  \partial_tv_*|_m=\frac{\langle(\dot v+v_P\,\partial_t\Ptop)/N_*^2\rangle}{\langle1/N_*^2\rangle},\qquad
  \int_0^{\Mtot}\partial_tv_*|_m\,\mathrm{d}m=0,
\end{equation}
the second because the box stays full, which fixes $\partial_t\Ptop$. In AN2D the levels are fixed and there is no lid:
$\partial_tv_*=\langle\dot v/N_*^2\rangle/\langle1/N_*^2\rangle$, $\partial_z(\partial_t\pi_*)=g\rho_R\partial_tv_*$,
$\partial_tP_*=\rho_R\partial_t\pi_*$. The weights $1/N_*^2$ differ between the parcels of a level only through $v_P$
(thermobaricity): a weakly stratified parcel moves further for a given change of volume and displaces more of its neighbours.
The plain mean, which a derivative at fixed placement gives, is exact only for a separable $v$.

\paragraph{Discrete form.} The parcels at a level are those of a slot, with the masses of the placement (entries of zero mass
skipped). For parcel $n$ in slot $k$, with $\dot v$ and $v_P$ at $(P_k,\vt_n,q_n)$,
\begin{equation}
  N^2_{nk}=g\rho_k\left(\partial_zv_*(z_k)+g\rho_kv_{P,nk}\right),\qquad w_{nk}=\frac{x_{nk}}{N^2_{nk}},\qquad
  \langle a\rangle_k=\frac{\sum_nw_{nk}a_{nk}}{\sum_nw_{nk}},
\end{equation}
with $\rho_k=\rho_R(z_k)$ (AN2D) or $1/v_{*,k}$ (FC2D) and $\partial_zv_*$ from the profile. A slot with $N^2_{nk}\le0$ for one
of its parcels takes the plain weights $w_{nk}=x_{nk}$ and is counted (\code{plain\_weight\_slots}).
\begin{itemize}
  \item \textbf{AN2D:} $\partial_tv_{*,k}=\langle\dot v\rangle_k$, and $\partial_t\pi_*$ by the rule of $\pi_*$: slopes
        $g\rho_R\partial_tv_{*,k}$, values by the trapezoidal rule in $P_R$, zero mean over the mass of the rows. Between the
        end rows and the walls only $\partial_tv_{*,e}$ is kept: $\partial_tb_e$ needs the mixed derivatives
        $\partial^2v/\partial P\partial\vt$, $\partial^2v/\partial P\partial q$, which the equation of state does not provide.
        The neglected term of $\partial_t\pi_*$, $-\tfrac12\partial_tb_e\Delta P^2$, is at most $2.6\times10^{-3}$ of
        $\max|\partial_t\pi_*|$ for $\gamma=3\times10^{-4}$\,Pa$^{-1}$ at $M_z=16$ (\code{tools/reference\_tendency\_study.jl}).
  \item \textbf{FC2D:}
        \begin{equation}
          \partial_t\Ptop=-\frac{\sum_kM_k\langle\dot v\rangle_k}{\sum_kM_k\langle v_P\rangle_k},\qquad
          \partial_tv_{*,k}=\langle\dot v\rangle_k+\langle v_P\rangle_k\,\partial_t\Ptop ,
        \end{equation}
        at fixed slot masses. The slot edges move at $\partial_tz_e=L_x^{-1}\sum_{k\text{ below }e}M_k\partial_tv_{*,k}$, zero at
        both walls, and their pressures at $\partial_t\Ptop$. $\partial_tP_*$ at fixed height is the derivative of the
        monotone cubic $P_*$ with respect to these motions of its nodes, by dual numbers (value $+\,\varepsilon\times$
        tendency for the heights, pressures and secants of the edges, through the same slopes and cubic as $P_*$); at an
        edge it is $\partial_t\Ptop+g\rho_*\partial_tz_e$.
\end{itemize}
$\int R_\Pi\,\mathrm{d}V$ (the exchange of the terms in $\partial_tP_*$ with the rest) vanishes to the resolution of the slots,
as the reference heights of the parcels of a mixed slot differ from their placement by up to a slot thickness.

% =====================================================================================
\section{Budgets of ACE and APE}\label{sec:budgets}
The budgets of $\rho\Pi_1$ and $\rho\Pi_2$ of the energy notes are diagnosed as the budget of $\rho P$ of the energetics is,
from the fluxes of the tendencies and the operators of the model (\code{src/reference\_state/budget.jl}). The derivatives
$\Sigma_\vt$, $\Sigma_q$ of the static energy are evaluated at three states of a parcel: the state itself $(z,p)$; the
intermediate state $(z,P_*(z))$, marked with a tilde ($\tilde\Sigma_\vt$); and the reference state $(\zs,P_*(\zs))$, marked
with a star ($\Sigma^*_\vt$). With
$\Delta_1\Sigma_\vt=\Sigma_\vt-\tilde\Sigma_\vt$ (zero in AN2D), $\Delta_2\Sigma_\vt=\tilde\Sigma_\vt-\Sigma^*_\vt$ (likewise for
$q$, by quadrature of $\partial_\vt v$, $\partial_qv$ in the pressure), at a fixed reference state
\begin{equation}
  \partial_t(\rho\Pi_i)=B_i\,\partial_t\rho+\Delta_i\Sigma_\vt\,\partial_t(\rho\vt)+\Delta_i\Sigma_q\,\partial_t(\rho q),
\end{equation}
with
\begin{equation}
  B_1=\Pi_1+(p-P_*)\Sigma_P-\vt\Delta_1\Sigma_\vt-q\Delta_1\Sigma_q,\qquad B_2=\Pi_2-\vt\Delta_2\Sigma_\vt-q\Delta_2\Sigma_q .
\end{equation}
In Bernoulli form, with the reversible flux $\mathbf J^{\rm rev}=\rho\mathbf u\,\Im(B)+\rho\vt\mathbf u\,\Im(\Delta\Sigma_\vt)+
\rho q\mathbf u\,\Im(\Delta\Sigma_q)$, the irreversible flux $\mathbf J^{\rm irr}=\Im(\Delta\Sigma_\vt)\mathbf J_\vt+
\Im(\Delta\Sigma_q)\mathbf J_q$, the conversion $C_B=-\rho\mathbf u\cdot\nabla B-\rho\vt\mathbf u\cdot\nabla\Delta\Sigma_\vt
-\rho q\mathbf u\cdot\nabla\Delta\Sigma_q$ to kinetic energy, and the heatings $Q_\vt=-\mathbf J_\vt\cdot\nabla\Sigma_\vt$,
$Q_q=-\mathbf J_q\cdot\nabla\Sigma_q$, $Q=\Sigma_\vt\sigma_\vt$ at each state, the product rule of the operators,
$\delta(\Im(a)b)=a\,\delta b+\Im(b\,\delta a)$, closes at each cell to round-off (in AN2D up to the divergence of $\rho_R\mathbf u$)
\begin{align}
  \partial_t(\rho\Pi_1)&=-\nabla\cdot(\mathbf J^{\rm rev}_{\rm ACE}+\mathbf J^{\rm irr}_{\rm ACE})-C_{\rm ACE}
     -(Q_\vt-\tilde Q_\vt)-(Q_q-\tilde Q_q)+Q-\tilde Q+R_{\Pi_1},\\
  \partial_t(\rho\Pi_2)&=-\nabla\cdot(\mathbf J^{\rm rev}_{\rm APE}+\mathbf J^{\rm irr}_{\rm APE})-C_{\rm APE}
     -(\tilde Q_\vt-Q^*_\vt)-(\tilde Q_q-Q^*_q)+\tilde Q-Q^*+R_{\Pi_2},
\end{align}
with the terms in $\partial_tP_*$ of the energy notes,
$R_{\Pi_1}=\rho\left(\Sigma_P(z,p)-\Sigma_P(z,P_*(z))\right)\partial_tP_*(z)$ (zero in AN2D) and
$R_{\Pi_2}=\rho\,\Sigma_P(z,P_*(z))\,\partial_tP_*(z)-\rho\,\Sigma_P(\zs,P_*(\zs))\,\partial_tP_*(\zs)$. The residuals of both
budgets are diagnostics. So are the continuous forms of the energy notes, $\mathbf J^{\rm rev}_{\rm ACE}=\rho\Pi_1\mathbf u+
\mathbf J_C^*$, $C_{\rm ACE}=C_{b,1}^*+C_D^*-\nabla\cdot\mathbf J_C^*$, $C_{\rm APE}=C_{b,2}^*$, with $\mathbf J_C^*=(p-P_*)\mathbf u$,
the buoyancy conversions $C_{b,1}^*=\rho g\rho_*(v-\tilde v)w$, $C_{b,2}^*=\rho g\rho_*(\tilde v-v_*)w$ ($\tilde v=v(P_*(z))$;
in AN2D $C_{b,1}^*=0$, $C_{b,2}^*=\rho_Rg\rho_R(v(P_R)-v_*)w$) and the pressure dilatation $C_D^*=(p-P_*)\nabla\cdot\mathbf u$ (zero in
AN2D), with the residuals of these manipulations, which converge with the grid. The conversion $C$ from $P$ to $K$ of the
energetics is $C_{\rm ACE}+C_{\rm APE}+C_{\rm bg}$, with the background conversion $C_{\rm bg}$ zero in the continuous equations.

% =====================================================================================
\section{Diagnostics}\label{sec:outputs}
The reference state is computed once per diagnostics session (\code{diagnostics(model; reference\_state)}), and its
diagnostics are entries of the book:
\begin{itemize}
  \item \textbf{On the cells:} $\zs$; ACE, APE, AE; the potentials at the intermediate and reference states and their
        differences; $B_1$, $B_2$; the fluxes, their divergences, conversions, heatings and buoyancy terms of the budgets;
        the diabatic rates; $\partial_tP_*(\zs)$; $R_{\Pi_1}$, $R_{\Pi_2}$; $\partial_t(\rho\Pi_1)$, $\partial_t(\rho\Pi_2)$ and the
        residuals.
  \item \textbf{Profiles on the $z$ axis:} $P_*$, $v_*$, $\partial_tP_*$.
  \item \textbf{Numbers:} $\int\rho\Pi_1$, $\int\rho\Pi_2$, $\int\rho\Pi$; the background potential energy; the certificate
        violations; the pivots, pricing passes and \code{sorted\_energy\_excess}; the parcels at a wall; the slots with
        plain weights; in FC2D $\Ptop$, its tendency, the volume residual and the iterations of the top pressure.
\end{itemize}

% =====================================================================================
\section{Verification and cost}\label{sec:tests}
The tests (\code{test/reference\_state.jl}, \code{transportation.jl}, \code{exact\_rearrangement.jl},
\code{reference\_tendency.jl}, \code{available\_energy\_budget.jl}) check, on small grids and for both models:
\begin{itemize}
  \item \textbf{Exact cases:} a state at rest is its own reference state (each parcel in its own row, $E_*=E_P$, $\Pi$ at
        round-off); a fully mixed fluid has $E_P-E_*$ and $\Pi$ at round-off (AN2D, uniform $\rho_R$) or $O(\Delta z^4)$; the
        analytic $E_P-E_*$ of a rearranged column with a separable equation of state; brute force over all assignments of
        8 parcels, separable and thermobaric.
  \item \textbf{The simplex:} on random cost matrices with unequal masses, feasibility and optimality checked independently
        of the solver; the linear assignment case against brute force; degenerate problems, including one on which a
        simplex without a strongly feasible tree can cycle; with $\gamma=0$ the exact method moves no mass; with strong
        thermobaricity it passes the certificate the sort fails and releases more energy.
  \item \textbf{Properties:} mass conservation of every parcel and slot; the column fills the box (FC2D); invariance under a
        periodic shift and under the origin of the geopotential; weak duality; $r(\zs)=0$ and $N_*^2\ge0$ at interior minima,
        $\Pi_1,\Pi_2\ge0$; the same results on serial and threaded managers; no state-sized allocations.
  \item \textbf{Tendency:} with $\gamma=0$ it matches centred differences of the reference state under $\pm\varepsilon$ times the
        diabatic rates; the tendencies of $\rho\Pi_1$, $\rho\Pi_2$ with the $R$ terms match differences with the reference
        state recomputed; without diffusion it vanishes; $\sum_kM_k\partial_tv_{*,k}=0$ (FC2D); $|\int R_\Pi|/\int|R_\Pi|$
        converges at first order ($2\times10^{-2}$ to $2\times10^{-3}$ from $M_z=8$ to $64$).
  \item \textbf{Weights under thermobaricity:} two density-compensated families mixed in every row, one heated: the
        $N_*^{-2}$-weighted mean is $26\%$, $8.5\%$ and $5.7\%$ from centred differences of the exact rearrangement at
        $K=20$, 40, 80, the plain mean $35\%$ (\code{tools/exact\_rearrangement\_study.jl} at full size).
  \item \textbf{Budgets:} the budgets of ACE and APE close at each cell to round-off, with and without the $R$ terms.
\end{itemize}

\paragraph{Measured} (\code{tools/exact\_rearrangement\_study.jl}, 2026-10-27 and 2026-10-31).
\begin{itemize}
  \item Thermohaline lock exchange ($\gamma=10^{-7}$\,Pa$^{-1}$, $64\times32$): the sort fails the certificate for at most 3
        nearly tied parcels, fixed by 2 or 3 pivots; the exact method changes $E_P-E_*$ by at most $9\times10^{-10}$ of it and
        $P_*$ by at most $2\times10^{-8}$\,Pa, with the same roughness in time. There the sort is the Lorenz state to
        round-off. \code{reference\_state!} takes 0.4--0.6\,ms.
  \item Every parcel misplaced by the sort (two compensated families, $\gamma=1.3\times10^{-4}$\,Pa$^{-1}$): 315, 1114 and
        7405 pivots in 21, 38 and 104 passes at $K=20$, 40, 80; 50\,ms at $K=80$, $N=1600$, of which pricing is two thirds.
  \item Float32 sees no thermobaric reordering for $\gamma$ up to $10^{-3}$\,Pa$^{-1}$, where Float64 makes 95 pivots that
        move mass: the reduced costs are below its tolerance. The exact method needs Float64.
\end{itemize}

% =====================================================================================
\section{Decisions}\label{sec:decisions}
Taken with the user, and governing the implementation:
\begin{itemize}
  \item \textbf{Exact solver:} a transportation simplex written in Julia, generic in the float type, started from the sorted
        placement at each snapshot (2026-10-27). The sort and the exact method stay separate methods, to be compared; the sort
        is the default, the examples use the exact one. The tree is updated in place after each pivot (2026-10-31).
  \item \textbf{Lorenz profile} from the slot-mean specific volumes, in both models; the duals certify the placement only
        (2026-10-27). The dual profile changes abruptly when the parcel split between two slots changes, and a split parcel has
        the same static energy at both slots along it, so its $\zs$ jumps between them; both made the budgets rough in time
        and space (the generation by two orders of magnitude in the lock exchange).
  \item \textbf{Reference height:} the global minimiser of $S$, not the nearest root of Tailleux (2013) nor the placement.
  \item \textbf{AN2D beyond the end rows:} $v_*$ follows the compression of the parcels of the end slot (2026-10-26);
        $\partial_tb_e$ is neglected, with its size measured (2026-10-27).
  \item \textbf{Tendency of the reference state} from the diabatic tendencies with $N_*^{-2}$ weights, valid for any equation of
        state.
  \item \textbf{FC2D $\partial_tP_*$ at fixed height} as the derivative of the monotone cubic with respect to its nodes, by dual
        numbers (2026-10-27). A hand-written tangent agreed to round-off at the same cost; applying
        $\partial_tP_*=\partial_t\Ptop+g\rho_*\partial_tz|_m$ between the edges with $\partial_tz|_m$ linear is not the derivative of
        the $P_*$ of $\Pi$ (off by 0.4 to 0.7\%).
  \item \textbf{FC2D reference heights in an exactly neutral layer} (2026-10-26, no change): identical parcels with $N_*^2=0$
        have an $S$ flat up to the wiggles of the monotone cubic, about $10^{-10}$\,J\,kg$^{-1}$ at $M_z=32$; the coarse search
        then takes the lowest tied level, the bottom wall. Their $\Pi_2$ is about $6\times10^{-10}$\,J\,kg$^{-1}$, and the ties
        vanish after the first time step.
\end{itemize}

% =====================================================================================
\appendix
\section{Not implemented, and future work}\label{app:future}

\subsection{Speedups of the exact rearrangement (deferred)}
\begin{itemize}
  \item \textbf{Banded pricing:} price the pairs within $|k-k_n|\le B$ slots of the sorted slot before the full pass. Pricing is
        two thirds of the time where every parcel is misplaced, so this is the next lever there.
  \item \textbf{Warm start from the previous snapshot.} The parcels are the cells, so after advection the previous placement has
        no direct relation to the new state; in AN2D it stays feasible (fixed masses), in FC2D the flows would need recomputing.
  \item \textbf{Warm start of the sort} from the previous order (insertion sort), with a fallback to the full sort when the flow
        overturns. At most about 1\,ms of 16\,ms per call at $256\times128$.
  \item \textbf{A starting interval for the Hermite search} in the refinement of $\zs$, which evaluates the profile in one
        bracket. Estimated 20--40\% of an evaluation; not measured.
\end{itemize}

\subsection{Other solvers}
\begin{itemize}
  \item \textbf{Greedy filling with repair.} Fill the slots from the top with the parcels of largest $v(P_k)$ at each slot
        pressure (Huang 2005; Wong et al.\ 2016; the discrete counterpart of Saenz et al.\ 2015), or bottom-up, or by divide and
        conquer at the pressure of the boundary that halves the mass (Stansifer et al.\ 2017), then repair by pairwise
        exchanges within a band of slots. No optimality guarantee: unrepaired heuristics can give $E_P-E_*<0$ (Harris \&
        Tailleux 2018), and 2-opt can stall (Hieronymus \& Nycander 2015). Not needed while the exact method is fast enough.
  \item \textbf{Histogram reduction} (Tseng \& Ferziger 2001; Saenz et al.\ 2015): bin the parcels in $(\vt,q)$ and solve for the
        bins. Averaging $(\vt,q)$ in a bin biases the energy under a nonlinear equation of state, and a resting stratified state no
        longer has zero available energy exactly; worthwhile only for much larger grids or long series.
  \item \textbf{Dual ascent:} maximise the concave piecewise-linear dual
        $D(\psi)=\sum_kM_k\psi_k+\sum_nm_n\min_k(c_{nk}-\psi_k)$ in $K$ variables, whose supergradient is the volume residual of each
        slot, by bundle or subgradient methods or entropic smoothing; the discrete counterpart of the semi-discrete optimal
        transport of Egan et al.\ (2025). Ties need a small flow among the tied slots.
  \item \textbf{Linear assignment} (Hungarian, Jonker--Volgenant) for equal masses with $K=N$: $O(N^3)$, and $10^3$--$10^5$ times
        slower than minimum-cost flow on the problem of Su \& Ingersoll (2016).
  \item \textbf{A guaranteed bound} (\S\ref{sec:complexity}): a polynomial pivot rule for the network simplex (Orlin 1997), or a
        strongly polynomial minimum-cost flow algorithm (Orlin 1993), with the same certificate. Likely slower on our states,
        where the sort is nearly optimal and the simplex needs 1 to 4 passes; worth it only if states appear on which the
        number of passes grows beyond about $K$.
\end{itemize}

\subsection{Other reference states}
\begin{itemize}
  \item \textbf{Horizontal-mean reference state:} $P_*$ hydrostatic with the row mean of the density (FC2D, by the model's
        discrete hydrostatic rule, anchored at one level) or of $v$ at $P_R$ (AN2D), with $\zs$ by the same LNB search. Not a
        rearrangement: $\int R_\Pi\,\mathrm{d}V\ne0$ and $\int\Pi\,\mathrm{d}m$ exceeds $E_P-E_*$ of the Lorenz state. Cheap, and
        local. Its tendency would follow from the model's full tendencies, as it changes under advection too. It was
        implemented with the first version of the Lorenz state and removed with the user (2026-10-26), as it is not a goal
        of the project.
  \item \textbf{Eddy-size constrained reference state} (Su \& Ingersoll 2016): forbid moves beyond a horizontal distance, which
        turns the problem back into a sparse assignment problem.
\end{itemize}

\subsection{Resolution and outputs}
\begin{itemize}
  \item \textbf{Slot resolution other than the rows} ($K\ne M_z$; $s$ slots per row in AN2D, equal-mass slots in FC2D), which
        the first design made a required input. With $K\ne M_z$ the actual state is no longer a placement with exactly the
        model's energy, and a state at rest gets $E_P-E_*=O(\Delta z^2)$.
  \item \textbf{Refinement of every near-lowest local minimum} of $S$, with a count of the parcels with two minima within a
        tolerance, where $\zs$ can jump.
  \item \textbf{Placement height} $\hat z_n=\sum_kx_{nk}z_k/m_n$ (or the sub-row height in sort order, the analogue of the
        Winters et al.\ height), as a check of $\zs$, and the reference static energy and $N_*^2$ per parcel as outputs.
  \item \textbf{$\partial_tP_*$ from differences of successive snapshots,} as a diagnostic. The tendency from the diabatic rates
        excludes the changes of the reference state caused by the discretisation of advection, and the differences include
        them, so their difference measures spurious mixing locally, the counterpart of the background potential energy of
        Il{\i}cak et al.\ (2012). Snapshot intervals must be long compared with $\Delta x/|\mathbf u|$, short compared with the
        diabatic evolution of $P_*$.
  \item \textbf{FC2D slot masses under advection:} the tendency holds the slot masses (the row masses of the state) fixed,
        although advection changes them.
  \item \textbf{$\partial_tb_e$ beyond the AN2D end rows,} once the equation of state provides the mixed second derivatives.
\end{itemize}

% =====================================================================================
\section*{References}
\begin{itemize}\small
  \item Ahuja, Magnanti \& Orlin (1993), \emph{Network Flows: Theory, Algorithms, and Applications}, Prentice Hall,
        ISBN 978-0-13-617549-0 (no DOI).
  \item Alefeld, Potra \& Shi (1995), Algorithm 748: enclosing zeros of continuous functions, ACM Trans.\ Math.\ Softw.\ 21,
        327--344, \doi{10.1145/210089.210111}.
  \item Bertsimas \& Tsitsiklis (1997), \emph{Introduction to Linear Optimization}, Athena Scientific, ISBN 978-1-886529-19-9
        (no DOI).
  \item Bradley, Brown \& Graves (1977), Design and implementation of large scale primal transshipment algorithms,
        Management Sci.\ 24, 1--34, \doi{10.1287/mnsc.24.1.1}.
  \item Cunningham (1976), A network simplex method, Math.\ Programming 11, 105--116, \doi{10.1007/BF01580379}.
  \item Dantzig (1951), Application of the simplex method to a transportation problem, in T.\ C.\ Koopmans (ed.),
        \emph{Activity Analysis of Production and Allocation}, Cowles Commission Monograph 13, Wiley, 359--373 (no DOI;
        \href{https://cowles.yale.edu/research/cfm-13-activity-analysis-production-and-allocation}{Cowles Foundation}).
  \item Egan, Methven, Bourne \& Cullen (2025), A definition of the background state of the atmosphere using optimal
        transport, arXiv:2504.11141.
  \item Glover, Klingman \& Stutz (1974), Augmented threaded index method for network optimization, INFOR 12, 293--298,
        \doi{10.1080/03155986.1974.11731583}.
  \item Grigoriadis (1986), An efficient implementation of the network simplex method, Math.\ Programming Study 26, 83--111,
        \doi{10.1007/BFb0121089}.
  \item Harris \& Tailleux (2018), Assessment of algorithms for computing moist available potential energy, QJRMS 144, 1501--1510.
  \item Hieronymus \& Nycander (2015), Finding the minimum potential energy state by adiabatic parcel rearrangements with a
        nonlinear equation of state: an exact solution in polynomial time, JPO 45, 1843--1857, \doi{10.1175/JPO-D-14-0174.1}.
  \item Hitchcock (1941), The distribution of a product from several sources to numerous localities, J.\ Math.\ Phys.\ 20,
        224--230, \doi{10.1002/sapm1941201224}.
  \item Hoffman (1963), On simple linear programming problems, in \emph{Convexity}, Proc.\ Symp.\ Pure Math.\ 7, AMS,
        317--327, \doi{10.1090/pspum/007/0157778}.
  \item Huang (2005), Available potential energy in the world's oceans, J.\ Mar.\ Res.\ 63, 141--158.
  \item Il{\i}cak, Adcroft, Griffies \& Hallberg (2012), Spurious dianeutral mixing and the role of momentum closure, Ocean
        Modelling 45--46, 37--58.
  \item Kleinschmidt \& Schannath (1995), A strongly polynomial algorithm for the transportation problem, Math.\ Programming
        68, 1--13, \doi{10.1007/BF01585755}.
  \item Orlin (1993), A faster strongly polynomial minimum cost flow algorithm, Oper.\ Res.\ 41, 338--350,
        \doi{10.1287/opre.41.2.338}.
  \item Orlin (1997), A polynomial time primal network simplex algorithm for minimum cost flows, Math.\ Programming 78,
        109--129, \doi{10.1007/BF02614365}.
  \item Saenz, Tailleux, Butler, Hughes \& Oliver (2015), Estimating Lorenz's reference state in an ocean with a nonlinear
        equation of state for seawater, JPO 45, 1242--1257.
  \item Stansifer, O'Gorman \& Holt (2017), Accurate computation of moist available potential energy with the Munkres
        algorithm, QJRMS 143, 288--292.
  \item Su \& Ingersoll (2016), On the minimum potential energy state and the eddy size-constrained APE density, JPO 46,
        2663--2674, \doi{10.1175/JPO-D-16-0074.1}.
  \item Tardos (1985), A strongly polynomial minimum cost circulation algorithm, Combinatorica 5, 247--255,
        \doi{10.1007/BF02579369}.
  \item Tailleux (2013), Available potential energy density for a multicomponent Boussinesq fluid with arbitrary nonlinear
        equation of state, JFM 735, 499--518.
  \item Tokuyama \& Nakano (1995), Efficient algorithms for the Hitchcock transportation problem, SIAM J.\ Comput.\ 24,
        563--578, \doi{10.1137/S0097539792236729}.
  \item Tseng \& Ferziger (2001), Mixing and available potential energy in stratified flows, Phys.\ Fluids 13, 1281--1293.
  \item Winters, Lombard, Riley \& D'Asaro (1995), Available potential energy and mixing in density-stratified fluids,
        JFM 289, 115--128, \doi{10.1017/S002211209500125X}.
  \item Winters \& D'Asaro (1996), Diascalar flux and the rate of fluid mixing, JFM 317, 179--193.
  \item Wong, Tailleux \& Gray (2016), The computation of reference state and APE production by diabatic processes in an
        idealized tropical cyclone, QJRMS 142, 2646--2657.
  \item Zadeh (1973), A bad network problem for the simplex method and other minimum cost flow algorithms, Math.\ Programming
        5, 255--266, \doi{10.1007/BF01580132}.
\end{itemize}

\end{document}
